% ----------------------------------------------------------
\chapter{Introdução}
\label{cap:introducao}
% ----------------------------------------------------------

O desenvolvimento de projetos geométricos rodoviários é uma etapa fundamental no processo de planejamento e implantação de infraestruturas de transporte, seguindo os princípios estabelecidos por \cite{lee2003}. A definição adequada do traçado de uma rodovia influencia diretamente o custo de construção, a segurança operacional e o impacto ambiental da obra. Assim, a fase inicial do projeto deve buscar a solução que concilie eficiência técnica, viabilidade econômica e respeito às condições naturais do terreno.

Este estudo tem como propósito o desenvolvimento da primeira etapa do projeto geométrico de um novo segmento rodoviário localizado no Distrito Federal, compreendendo o traçado entre dois pontos obrigatórios de coordenadas e cotas conhecidas — o Ponto de Partida (PP) e o Ponto Final (PF). A concepção geométrica foi elaborada conforme as diretrizes da norma vigente para rodovias de Classe I-B \cite{dner1999}, adotando como veículo de projeto o tipo "CO".


A análise foi conduzida a partir de dados topográficos obtidos na plataforma Geoportal DF (SICAD), que serviram como base para a modelagem digital do terreno. As superfícies, perfis longitudinais e seções transversais foram processados utilizando o software Autodesk Civil 3D, o que possibilitou a obtenção de parâmetros geométricos, altimétricos e planimétricos necessários à definição dos alinhamentos.

Foram desenvolvidas duas alternativas de traçado, denominadas Traçado A e Traçado B, as quais foram avaliadas segundo critérios técnicos e ambientais, buscando a alternativa mais adequada para implantação. A comparação entre os traçados considerou aspectos como declividade, comprimento total, volume de movimentação de terras e adequação ao relevo natural, além de potenciais interferências ambientais e de ocupação do solo.

O presente relatório apresenta as etapas de desenvolvimento e análise do estudo de traçado, abrangendo desde a caracterização topográfica da área até o detalhamento do alinhamento selecionado. Os resultados obtidos buscam subsidiar a escolha da solução de engenharia mais eficiente, conciliando economia de implantação, funcionalidade e sustentabilidade.
